% Frequency scaling of NLOS path loss in a ray-traced street canyon.
% IEEE conference format (IEEEtran). Compiles on Overleaf with no extra setup.
%
% STATUS: complete draft, all sections. The abstract is sized for a standard
% six-page conference paper and will need trimming if the venue is shorter.
%
% FIGURES NEEDED (generated by the repository scripts, not committed).
% Both are double-column figure*, so render with --ieee, which sizes them to
% 7.16 in and sets the fonts for that width:
%   figures/umi_comparison_uma_ieee.png
%     python scripts/15_umi_figure.py --data data/bands_los.npz \
%            --outdir paper/figures --ieee --model uma
%   figures/fig_roughness_uma_ieee.png
%     python scripts/19_roughness_analysis.py --model uma \
%            --out results/roughness_summary_uma.json
%     python scripts/20_roughness_figure.py --model uma \
%            --summary results/roughness_summary_uma.json \
%            --out paper/figures/fig_roughness_uma_ieee.png --ieee
% Drop them in paper/figures/ before compiling, or comment out the figure blocks.

\documentclass[conference]{IEEEtran}
\IEEEoverridecommandlockouts

\usepackage{cite}
\usepackage{amsmath,amssymb}
\usepackage{graphicx}
\usepackage{booktabs}
\usepackage{siunitx}
\usepackage{url}

% Float parameters.  Five tables and two double-column figures is more than
% LaTeX's defaults will place: topnumber 2 and totalnumber 3 starve the
% figures and defer them past the references.  Raising the limits and the
% fractions keeps each float near the text that discusses it.
\setcounter{topnumber}{3}
\setcounter{totalnumber}{4}
\renewcommand{\topfraction}{0.92}
\renewcommand{\textfraction}{0.07}
\renewcommand{\floatpagefraction}{0.75}

\sisetup{detect-all, per-mode=symbol}
\DeclareSIUnit{\dBm}{dBm}
\graphicspath{{figures/}}

\begin{document}

\title{Frequency Scaling of Non-Line-of-Sight Path Loss in a\\
Ray-Traced Urban Street Canyon}

% PROPOSED author list, not an agreed one. Ali told Thompson and Qi on
% 11 Oct 2026 that nothing goes out with their names on it unless they agree,
% so the covering email must say plainly that this block is a proposal and
% that the single-author version is one edit away.
%
% Affiliation lines are kept SHORT on purpose. IEEEtran lays three author
% blocks out side by side, roughly 2.4 in each, and the full strings
% ("School of Information Science and Engineering", "Institute for Imaging,
% Data and Communications") overflow into the neighbouring block. If they
% want the full wording, break it across more \\ lines rather than widening.
%
% TO CONFIRM WITH THEM: author order, how they want their names written,
% Qi's affiliation (visiting Edinburgh from Southeast University, so it may
% be one or both), and whether they want their email addresses printed.
%
% IF THEY DECLINE: delete the second and third blocks here and uncomment the
% Acknowledgements section at the end of the file.
\author{
\IEEEauthorblockN{A. J. Taghiyev}
\IEEEauthorblockA{\textit{Kings' School Al Barsha}\\
Dubai, United Arab Emirates}
\and
\IEEEauthorblockN{T. Qi}
\IEEEauthorblockA{\textit{Southeast University}\\
Nanjing, China}
\and
\IEEEauthorblockN{J. S. Thompson}
\IEEEauthorblockA{\textit{IDCOM, School of Engineering}\\
\textit{University of Edinburgh}\\
Edinburgh, United Kingdom}
}

\maketitle

\begin{abstract}
A coverage model validated at one carrier can still be wrong about the next,
and a single-band residual cannot show it. We measure the quantity that can:
path loss at the same receiver location across several carriers. One urban
street canyon is traced with Sionna RT at \SI{3.5}{}, \SI{28}{} and
\SI{39}{\giga\hertz}, and a frequency slope in decibels of path loss per
decibel of carrier frequency is estimated for each cell from its three
simulated values, then compared against 3GPP TR 38.901. The transmitter is a
rooftop site, so the urban macrocell model applies rather than the microcell
one; against it, line-of-sight cells agree to within \SI{0.3}{\decibel} with
\SI{99}{\percent} interval coverage.
Non-line-of-sight cells do not, and the disagreement is principally about the
slope. Beyond \SI{80}{\metre}, \SI{81.0}{\percent} of non-line-of-sight
cells have slopes between \num{2.9} and \num{3.1}, against a reference value
of \num{2.00}; a slope of \num{3.0} is what power arriving over a single
diffracting edge gives. Comparing slopes cell by cell at two interaction
depths leaves the far-zone median unchanged, so this follows from which
mechanisms are enabled rather than from how far paths are traced. Giving every
surface a scattering coefficient derived from its roughness rather than fitted
carries those far-zone slopes down through the reference value between
\SI{0.2}{} and \SI{0.5}{\milli\metre} of RMS roughness, while the fraction
of cells inside the model's shadow fading interval improves separately and
monotonically, saturating near \SI{52}{\percent}. We also report that the ray
budget
needed to converge a map does not vary with carrier frequency, but that
enabling diffuse scattering raises it by orders of magnitude and leaves the
slopes above unconverged at the budgets used.
\end{abstract}

\begin{IEEEkeywords}
non-line-of-sight path loss, frequency scaling, ray tracing, diffuse
scattering
\end{IEEEkeywords}

\section{Introduction}
\label{sec:intro}

A coverage model can be validated at one carrier and still be wrong about the
next. The question this paper asks is narrow and, as far as we can establish,
unmeasured: at a fixed receiver location in deep shadow, how does path loss
change as the carrier moves from sub-6 to millimetre wave, and does a
site-specific ray tracer change it the same way a standardised model says it
should?

A single-band check cannot answer this. Agreement at \SI{3.5}{\giga\hertz}
constrains the level of a map, not its frequency derivative, and two
predictors that match closely at one carrier are free to diverge at the next
in a way no single-band residual exposes. The quantity that would expose it is
the paired one: path loss at the same receiver location, at several carriers,
differenced against itself rather than aggregated over cells.

What the literature establishes falls either side of that quantity. The 38.901
urban models \cite{3gpp38901} are fitted to multiband measurement campaigns
\cite{rappaport2013mmwave, sulyman2014pathloss}, which report path loss
exponents and shadow fading per band and per environment class. Ray tracers
are calibrated against measurement and reported in the same terms
\cite{ying2026nyuray}. Both therefore validate aggregate fits over a
population of locations.
% TODO (Ali): Qi asked specifically what the closest multiband studies
% validated and whether any reports the spatial distribution of paired
% per-cell frequency slopes. The sentence above says only what the fitted
% quantities are, which is safe. If you read [2] and [3] closely enough to
% say that none reports a per-location slope distribution, say it outright,
% because it is the gap the paper claims. Do not assert it until you have.
Diffuse scattering has its own modelling literature
\cite{degliesposti2001diffuse, degliesposti2004advanced}, with measured
coefficients for real facades \cite{degliesposti2007measurement}, but again
characterised as a surface property rather than through its effect on
cross-band behaviour at a location. Meanwhile traced maps increasingly serve
as training data and reference for learned radio map estimators
\cite{feng2025radiomapsurvey, yapar2024challenge}, and the solver settings
behind those datasets are chosen for tractability: the benchmark much of that
literature trains on caps paths at two interactions and reports no ray count
\cite{yapar2024challenge}. What is missing in all of this is the per-location
frequency slope itself.

A simulation can supply it, because the same geometry can be traced at several
carriers and the cells put in correspondence exactly. We trace one urban
street canyon at \SI{3.5}{}, \SI{28}{} and \SI{39}{\giga\hertz} with Sionna RT
\cite{hoydis2023sionnart} and estimate a frequency slope for each receiver
cell from its three simulated path-loss values, by regression on
$10\log_{10}(f_c)$. The slope is dimensionless: decibels of path loss per
decibel of carrier frequency. Free space gives \num{2.0}.

Three results follow.

First, whether the simulator and the model agree depends entirely on
visibility, and on which reference model the geometry calls for. The
transmitter is a rooftop site, so the macrocell model applies rather than the
microcell one despite the scene being a canyon; against it, line-of-sight
cells agree to within \SI{0.3}{\decibel} with \SI{99}{\percent} interval
coverage, while non-line-of-sight cells carry a bias of \SI{9}{} to
\SI{16}{\decibel} that is not a constant offset but two regimes, reflection
lit near the transmitter and diffraction lit beyond it.

Second, the simulated frequency slope in deep shadow is larger than the
reference, by an amount with a clean physical reading. Beyond \SI{80}{\metre},
\SI{81.0}{\percent} of non-line-of-sight cells have slopes between \num{2.9}
and \num{3.1}, against the macrocell model's \num{2.00}. A slope of \num{3.0}
is what power arriving over a single diffracting edge gives, since the uniform
theory of diffraction \cite{kouyoumjian1974utd} makes the diffraction
coefficient fall as $1/\sqrt{k}$, where $k = 2\pi/\lambda$ is the
wavenumber. Comparing slopes cell by cell at two
interaction depths shows the median paired change is zero, so this follows
from which mechanisms are enabled rather than from how far paths are traced.

Third, giving every surface a scattering coefficient derived from its
roughness rather than fitted moves those far-zone slopes down through the
reference value as roughness rises from \SI{0.2}{} to
\SI{0.5}{\milli\metre}, while the fraction of cells inside the model's
shadow fading interval improves separately and monotonically. The two measures
do not peak at the same roughness, and the scattering-on slopes are not
converged at the budgets used, so what the sweep establishes is the direction
and scale of the effect rather than a crossing point. It is enough to identify
the mechanism: a tracer run without diffuse scattering lights the deep shadow
with one diffracting edge and inherits that edge's frequency dependence.

Section~\ref{sec:setup} gives the scene, the reference models and the
convergence of the traced maps. Section~\ref{sec:umi} compares the maps
against 38.901 band by band. Section~\ref{sec:freq} measures the per-cell
frequency slope and tests the diffraction reading. Section~\ref{sec:scattering}
introduces roughness-derived diffuse scattering and follows the far-zone
slopes through it.

% =====================================================================
\section{Simulation Setup and Reference Models}
\label{sec:setup}

\subsection{Scene and solver configuration}

All results use the \texttt{simple\_\hspace{0pt}street\_\hspace{0pt}canyon} scene shipped with Sionna
RT 2.0.1 (Mitsuba 3.8.0, Dr.Jit 1.3.1, CPU LLVM backend) and a single
isotropic transmitter at $(-33, 11, 32)$\,m radiating
\SI{30}{\dBm} with vertical polarisation. The scene contains four buildings
with roof heights of \SI{21.8}{}, \SI{29.1}{}, \SI{50.9}{} and
\SI{50.9}{\metre}, and the transmitter stands on the footprint of one of the
lower pair, so it is a rooftop site above two of the four roofs and below the
other two. This geometry matters for the choice of reference model and is
returned to in Section~\ref{sec:umi}. The radio map is evaluated on a
horizontal plane at \SI{1.5}{\metre} with \SI{1}{\metre} square cells, giving a
$122 \times 187$ grid of \num{22814} cells. Propagation is traced to a maximum
depth of three, with line of sight, specular reflection and diffraction
enabled. Diffuse scattering is disabled for the baseline of
Sections~\ref{sec:umi} and~\ref{sec:freq} and is the subject of
Section~\ref{sec:scattering}.

Two configuration choices are worth stating because they change what the map
means. First, refraction is enabled by default in Sionna and is switched off
here: left on, rays pass through building walls, and cells inside buildings
stop being structurally unreachable. Second, cells lying under building
geometry are identified by casting a vertical ray upward from each cell centre
rather than by thresholding received power, which would discard genuine
deep-shadow street cells. This removes \num{5236} cells.

Line of sight is decided geometrically, by a single shadow ray from each cell
centre to the transmitter, so the classification is exact and independent of
the Monte Carlo seed.

Every statistic below is computed on one common mask, fixed once and used for
all three bands and all solver settings, so that changes in coverage are never
confused with changes in precision on the same cells. Of the \num{22814} grid
cells, \num{5236} lie under building geometry and \num{17578} are street.
Of those, \num{202} are closer than \SI{10}{\metre} in two dimensions and are
excluded, since the reference models of Section~\ref{sec:umi} are specified
only from \SI{10}{\metre}, and a further \num{31} receive no ray in at least
one band at \num{2e8} rays. Those \num{31} all lie between \SI{97}{} and
\SI{101}{\metre}, at the far edge of the map. The common mask is therefore
\num{17345} cells, \num{3285} line of sight and \num{14060} non line of
sight. Path loss for a cell is the negative of the solver's returned path
gain, which is an area average over the \SI{1}{\metre} cell, while the
line-of-sight classification uses the cell centre.

\subsection{Which reference model the geometry calls for}

38.901 specifies its urban models at fixed base station heights: Table 7.4.1-1
gives \SI{10}{\metre} for the urban microcell (UMi) street canyon and
\SI{25}{\metre} for urban macrocell (UMa). The transmitter here is a rooftop
site at \SI{32}{\metre}, above two of the four roofs in the scene, so neither
height matches exactly and UMa is much the closer of the two.

The scene's name invites UMi, and the layout is a street canyon, but the
antenna is not where UMi puts it. A matched-height UMi run is not available in
this scene: the transmitter stands on a building footprint, so lowering it to
\SI{10}{\metre} places it inside the building. Traced at that height it
reaches \num{988} of \num{22814} cells and no cell has line of sight.

We therefore take UMa as the primary reference and report UMi alongside it,
because UMi is what the scene description alone would suggest and the gap
between the two is itself informative. Throughout, the residual is defined as
ray-traced path loss minus model path loss, so a positive residual means the
ray tracer loses more power than the model predicts.

\subsection{Convergence of the ray-traced path-loss maps}

A ray-traced radio map is a Monte Carlo estimate, and any claim made by
differencing it against a model is only meaningful if the estimator noise is
small relative to the difference. This is rarely checked. We therefore trace
the identical scene several times and measure how far the maps differ from one
another.

Two properties of the solver have to be recorded first. Sionna RT 2.0.1 on the
LLVM backend does not reproduce a radio map from its seed: at \num{5e6} rays,
eight runs at one seed produced at least three distinct outcomes differing by
approximately \SI{0.8}{\decibel}. The behaviour is intermittent and shrinks
rapidly with ray count, with no pair of four runs differing at \num{2e7} rays,
so it is negligible at the operating point used here. The quantity reported
below is accordingly labelled run-to-run rather than seed-to-seed.

Table~\ref{tab:convergence} gives the run-to-run RMSE at \SI{3.5}{\giga\hertz}
over street cells reached in every run, together with the fraction of street
cells no ray reaches. The range spanned is wide, and the low end is not
exotic: Sionna's own default of \num{1e6} rays per transmitter leaves
\SI{9.1}{\percent} of street cells unreached and the rest disagreeing between
runs by \SI{4.08}{\decibel}. Three orders of magnitude further up, at
\num{2e8}, the unreached fraction has fallen to \SI{0.18}{\percent} and the
disagreement to \SI{0.32}{\decibel}. From \num{1e6} rays upward the error falls
as $N^{-1/2}$, as an unbiased Monte Carlo estimator should.

The practical consequence is that the ray budget sets a floor on what any
quantity computed by differencing such a map can resolve. Below roughly
\num{1e7} rays that floor is several decibels, which is the same order as the
discrepancies reported in Sections~\ref{sec:umi} and~\ref{sec:freq}. The budget
therefore has to be established before the comparison, not assumed.

\begin{table}[t]
\caption{Run-to-run agreement of the ray-traced map at \SI{3.5}{\giga\hertz}.
Solve times are the mean of the runs at each budget, measured on two threads
of an Intel Xeon at \SI{2.1}{\giga\hertz} with Mitsuba's
\texttt{llvm\_\hspace{0pt}ad\_\hspace{0pt}mono\_\hspace{0pt}polarized}
variant, timing the solver call only and excluding scene loading and
first-call compilation.}
\label{tab:convergence}
\centering
\begin{tabular}{rrrr}
\toprule
Rays per Tx & Street cells & Run-to-run & Solve \\
            & with no ray  & RMSE (dB)  & time (s) \\
\midrule
\num{2e5} & \SI{28.94}{\percent} & 4.75 & $<1$ \\
\num{1e6} & \SI{9.10}{\percent}  & 4.08 & 1 \\
\num{5e6} & \SI{0.98}{\percent}  & 2.49 & 3 \\
\num{2e7} & \SI{0.20}{\percent}  & 1.15 & 9 \\
\num{5e7} & \SI{0.18}{\percent}  & 0.67 & 22 \\
\num{1e8} & \SI{0.18}{\percent}  & 0.46 & 44 \\
\textbf{\num{2e8}} & \textbf{\SI{0.18}{\percent}} & \textbf{0.32} & \textbf{86} \\
\num{5e8} & \SI{0.18}{\percent}  & 0.20 & 269 \\
\bottomrule
\end{tabular}
\end{table}

Because this study compares three carrier frequencies, the ray budget was
measured separately in each rather than assumed to carry over. It does not move
with frequency. Fitting the run-to-run RMSE as $N^{\alpha}$ over \num{5e6} to
\num{5e8} rays gives $\alpha = -0.566$, $-0.570$ and $-0.575$ at
\SI{3.5}{}, \SI{28}{} and \SI{39}{\giga\hertz} respectively, and the ray count
required to place a single map \SI{0.25}{\decibel} from the converged mean is
\num{1.69e8}, \num{1.74e8} and \num{1.74e8}. The spread is under
\SI{3}{\percent} and runs marginally against the millimetre bands rather than
in their favour.

The reason is geometric. The variance in a cell is set by how many ray tubes
reach it, and frequency changes what each contribution is worth in decibels
rather than how many arrive. This was checked directly rather than inferred:
two runs at different frequencies disagree about which cells received a ray no
more than two runs at the same frequency do, and the fraction of unreached
street cells matches across bands to \SI{0.2}{} percentage points at every ray
count.

All maps in Sections~\ref{sec:umi} and~\ref{sec:freq} are therefore traced at
\num{2e8} rays per transmitter, where a single map sits
\SI{0.23}{\decibel} from the converged mean, roughly two orders of magnitude
below the discrepancies reported in Section~\ref{sec:umi}.

% =====================================================================
\section{Comparison with the 3GPP 38.901 Reference Models}
\label{sec:umi}

The traced maps are now compared cell by cell against both reference models at
all three carriers, on the common mask of Section~\ref{sec:setup}.


\begin{table*}[t]
\caption{Ray-traced path loss against the 38.901 UMa and UMi models.
Residual is ray tracer minus model. ``Residual SD'' is the spread of the
residual; the interval-coverage column counts cells whose residual falls
inside the model's own shadow fading standard deviation, which is
\SI{4.0}{\decibel} for LOS in both models, \SI{6.0}{\decibel} for UMa NLOS
and \SI{7.82}{\decibel} for UMi NLOS. The distance exponent is fitted to the
traced cells, with the model's own value in parentheses; it is a fit against
distance and is not the per-cell frequency slope of
Section~\ref{sec:freq}.}
\label{tab:umi}
\centering
\begin{tabular}{llrrrrrrrr}
\toprule
& & \multicolumn{4}{c}{UMa (\SI{25}{\metre} nominal)} & \multicolumn{4}{c}{UMi (\SI{10}{\metre} nominal)} \\
\cmidrule(lr){3-6}\cmidrule(lr){7-10}
Band & Class & Bias & Res.\ SD & Within & Dist. & Bias & Res.\ SD & Within & Dist. \\
     &       & (dB) & (dB) & $\pm 1\sigma$ & exp. & (dB) & (dB) & $\pm 1\sigma$ & exp. \\
\midrule
\SI{3.5}{\giga\hertz} & LOS  & $-0.09$ & 1.10 & \SI{99.5}{\percent} & 1.83 (2.20) & $-2.69$ & 1.04 & \SI{88.8}{\percent} & 1.83 (2.10) \\
                      & NLOS & $+9.43$ & 11.67 & \SI{15.0}{\percent} & 6.53 (3.91) & $+6.96$ & 11.90 & \SI{20.7}{\percent} & 6.53 (3.53) \\
\SI{28}{\giga\hertz}  & LOS  & $+0.26$ & 0.99 & \SI{99.4}{\percent} & 1.85 (2.20) & $-2.34$ & 0.92 & \SI{90.8}{\percent} & 1.85 (2.10) \\
                      & NLOS & $+14.72$ & 14.93 & \SI{15.2}{\percent} & 7.88 (3.91) & $+11.08$ & 15.19 & \SI{16.7}{\percent} & 7.88 (3.53) \\
\SI{39}{\giga\hertz}  & LOS  & $+0.26$ & 1.00 & \SI{99.3}{\percent} & 1.85 (2.20) & $-2.33$ & 0.93 & \SI{90.9}{\percent} & 1.85 (2.10) \\
                      & NLOS & $+15.60$ & 15.51 & \SI{14.9}{\percent} & 8.08 (3.91) & $+11.78$ & 15.77 & \SI{17.2}{\percent} & 8.08 (3.53) \\
\bottomrule
\end{tabular}
\end{table*}

Table~\ref{tab:umi} shows the two classes behaving quite differently, and the
choice of reference model separating them further.

Fig.~\ref{fig:umi}(a) shows the two classes against distance at
\SI{28}{\giga\hertz}. In line of sight the traced map reproduces UMa almost
exactly. The bias is
under \SI{0.3}{\decibel} in every band and \SI{99.3}{} to \SI{99.5}{\percent}
of cells fall inside the model's own one-sigma interval, against the
\SI{68}{\percent} that interval nominally covers. Against UMi the same cells
carry a negative bias of \SI{2.3}{} to \SI{2.7}{\decibel}, which is the
difference the \SI{22}{\metre} of antenna height between the two model
geometries accounts for. The residual spread is about \SI{1}{\decibel} either
way, so the LOS disagreement is a constant offset set by the reference model
rather than anything the tracer is doing.

Non-line-of-sight cells do not follow. Against UMa the bias runs from
\SI{9.4}{\decibel} at \SI{3.5}{\giga\hertz} to \SI{15.6}{\decibel} at
\SI{39}{\giga\hertz}, and interval coverage sits near \SI{15}{\percent} in
every band. Against UMi the bias is smaller but still grows the same way, from
\SI{7.0}{} to \SI{11.8}{\decibel}. Switching reference model moves the NLOS
numbers without changing their shape: the gap grows with frequency under both,
by \SI{6.2}{\decibel} across the three bands under UMa and \SI{4.8}{\decibel}
under UMi. Whatever is happening is not an artefact of picking the wrong
model.

Fig.~\ref{fig:umi}(b) shows why the NLOS numbers in Table~\ref{tab:umi}
should be read with care. The fitted distance exponents of \num{6.5} to
\num{8.1} should not be read as path loss exponents. A single straight line in $\log_{10} d$ is being
forced through two regimes that behave differently. Within roughly
\SI{45}{\metre} of the transmitter the ray tracer loses up to
\SI{10}{\decibel} \emph{less} than UMi, because those NLOS cells are
illuminated by specular reflection from the canyon walls. Beyond
\SI{80}{\metre} it loses \SI{13}{} to \SI{25}{\decibel} \emph{more}, and that
excess is larger at \SI{28}{} and \SI{39}{\giga\hertz} than at
\SI{3.5}{\giga\hertz}. The NLOS gap is therefore not a single offset that could
be absorbed into a constant, and the way it grows with frequency is the subject
of the next section.

\begin{figure*}[!t]
\centering
\includegraphics[width=\textwidth]{umi_comparison_uma_ieee.png}
\caption{Ray-traced street canyon against 3GPP 38.901 UMa at \SI{3.5}{},
\SI{28}{} and \SI{39}{\giga\hertz}. (a) Path loss against 3D distance at
\SI{28}{\giga\hertz}: dots are traced cells, lines are the UMa LOS and NLOS
means, shaded bands are the model's own $\pm 1\sigma$ shadow fading. (b) NLOS
residual against distance, median and interquartile range per band, with the
model's \SI{6}{\decibel} NLOS shadow fading in grey. (c) Distribution of the
per-cell frequency coefficient. UMa's NLOS frequency term is
$20\log_{10}(f_c)$, so its reference coincides with free space; the two peaks
sit on that value and on the single-edge value of \num{3.0}.
(d) The same coefficient mapped over the scene; grey is building footprint,
the star is the transmitter, and the dashed arc is the \SI{80}{\metre}
boundary used to define the far zone. The colour scale is clipped to
\numrange{1.8}{3.1}, which spans all but the lowest percentiles; cells below
\num{1.8} take the darkest colour and cells reached by no ray are left as
building grey.}
\label{fig:umi}
\end{figure*}

% =====================================================================
\section{How Path Loss Scales with Frequency}
\label{sec:freq}

Because the same geometry is traced at three carriers, the frequency dependence
can be measured per cell rather than inferred from aggregate fits. Regressing
each cell's path loss on $10\log_{10}(f)$ across the three bands yields one
coefficient per cell. Free space gives \num{2.0}; the 38.901 UMi NLOS formula
gives \num{2.13}, and the 38.901 UMa NLOS formula gives exactly \num{2.00},
since its only frequency term is $20\log_{10}(f_c)$. Against the macrocell
model that this geometry calls for, the reference and free space coincide.

\begin{table}[t]
\caption{Per-cell frequency coefficient, from regressing path loss on
$10\log_{10}(f)$ across the three carriers.}
\label{tab:freqcoef}
\centering
\begin{tabular}{lrrrr}
\toprule
 & Median & 10th--90th & 3.5--28 & 28--39 \\
 &        & percentile & \si{\giga\hertz} & \si{\giga\hertz} \\
\midrule
LOS  & 2.03 & 2.00--2.07 & 2.03 & 2.00 \\
NLOS & 2.99 & 1.76--3.00 & 2.99 & 3.00 \\
\bottomrule
\end{tabular}
\end{table}

Fig.~\ref{fig:umi}(c) gives the distribution. Line-of-sight cells sit on free
space, as the model assumes. Non-line-of-sight cells do not: the median coefficient is \num{2.99}, against the model's
\num{2.00} under UMa and \num{2.13} under UMi. This is not a broad
distribution with a misleading centre. The
NLOS cells divide as \SI{57.5}{\percent} within \num{0.1} of \num{3.0},
\SI{22.9}{\percent} between \num{2.1} and \num{2.9}, \SI{6.4}{\percent}
within \num{0.1} of free space, and \SI{13.2}{\percent} below \num{1.9}; no
cell exceeds \num{3.1}. Fig.~\ref{fig:umi}(d) maps the same quantity over the scene, and the share
near \num{3.0} rises sharply with distance from the transmitter, from \SI{19.8}{\percent} within \SI{40}{\metre}
($n = \num{1855}$), to \SI{50.2}{\percent} between \SI{40}{} and
\SI{80}{\metre} ($n = \num{7017}$), to \SI{81.0}{\percent} beyond
\SI{80}{\metre} ($n = \num{5188}$), where the median is \num{2.99}.

A coefficient of exactly \num{3.0} has a clean physical reading. The uniform
theory of diffraction gives a diffraction coefficient falling as
$1/\sqrt{k}$ in the wavenumber $k = 2\pi/\lambda$, so a single diffracting
edge contributes
\SI{10}{\decibel} per decade of frequency on top of the \SI{20}{\decibel} of
free space. A cell whose received power arrives over one edge therefore scales
as $30\log_{10}(f)$. The \SI{19.6}{\percent} of NLOS cells sitting at
\num{2.1} or below are the reflection-lit set identified in
Section~\ref{sec:umi}; their small departures from \num{2.0} are consistent
with the ITU material parameters \cite{iturp2040} varying with frequency. The
\SI{22.9}{\percent} between the two groups are reached by both mechanisms and
lie mostly in the \SI{40}{} to \SI{80}{\metre} band, where the two regimes
overlap. Integrated over \SI{3.5}{} to \SI{39}{\giga\hertz}, a cell scaling at
\num{3.0} rather than the macrocell model's \num{2.00} loses
\SI{10.5}{\decibel} more; against UMi's \num{2.13} the figure is
\SI{9.1}{\decibel}. The excess over the macrocell reference is almost exactly
the \SI{10}{\decibel} per decade that one diffracting edge contributes, which
is what makes the reading unambiguous.

This matters beyond the comparison itself. Learned path loss and radio map
models take frequency as an input \cite{nguyen2021multifreq} and, when asked
to generalise to an unseen band, typically do so by predicting a residual on
top of an analytic prior, on the assumption that the prior carries most of the
frequency dependence \cite{hussain2026zeroshot}. A free-space prior carries a
coefficient of \num{2.0}, so that assumption can be tested directly here: take
the \SI{3.5}{\giga\hertz} map, add $20\log_{10}(f/3.5)$, and score it against
the traced map at the other two carriers on the same cells. In line of sight
the translated map is accurate to \SI{0.43}{\decibel} RMSE at
\SI{28}{\giga\hertz} and \SI{0.44}{\decibel} at \SI{39}{\giga\hertz}. In non
line of sight it is wrong by \SI{7.41}{} and \SI{8.44}{\decibel}, biased by
\SI{5.29}{} and \SI{6.17}{\decibel} towards predicting too little loss. A
free-space prior therefore transfers almost perfectly where the path is
unobstructed and fails by several decibels exactly where coverage prediction
is hardest.

\subsection{Effect of maximum interaction depth}

The obvious alternative explanation is that the solver simply stops too early,
and that deeper multiple-bounce paths would carry the missing energy into the
shadow. This was tested directly by retracing all three bands at a maximum
depth of five with the same ray count and seed, and comparing the per-cell
frequency slope between the two depths on the same cells.

\begin{table}[t]
\caption{Effect of raising maximum interaction depth from three to five.
Biases are against UMa. The lower block reports the per-cell frequency slope,
which the band-wise biases above do not test directly.}
\label{tab:depth}
\centering
\begin{tabular}{lrrrr}
\toprule
Band & \multicolumn{2}{c}{NLOS bias (dB)} & \multicolumn{2}{c}{Beyond \SI{80}{\metre} (dB)} \\
\cmidrule(lr){2-3}\cmidrule(lr){4-5}
 & depth 3 & depth 5 & depth 3 & depth 5 \\
\midrule
\SI{3.5}{\giga\hertz} & $+9.43$ & $+9.27$ & $+13.92$ & $+13.74$ \\
\SI{28}{\giga\hertz}  & $+14.72$ & $+13.90$ & $+21.56$ & $+20.40$ \\
\SI{39}{\giga\hertz}  & $+15.60$ & $+14.58$ & $+22.85$ & $+21.44$ \\
\midrule
\multicolumn{5}{l}{\emph{Median per-cell frequency slope}} \\
all NLOS        & 2.99 & 2.76 & & \\
$>\SI{80}{\metre}$  & 2.99 & 2.99 & & \\
\multicolumn{3}{l}{median paired change $0.000$} & & \\
\multicolumn{3}{l}{cells moving by $>0.05$: \SI{11.1}{\percent}} & & \\
\bottomrule
\end{tabular}
\end{table}

It is not, and the per-cell slope is the test that settles it. No cell is
reached at depth five that was not already reached at depth three, and the
median NLOS path loss does not move at any of the three carriers. The mean
NLOS bias falls by \SI{0.16}{} to \SI{1.02}{\decibel}, which on its own says
little about the frequency slope, since two band-specific biases can change
without the slope between them changing at all.

Comparing the slopes cell by cell is unambiguous. The median paired change is
zero to four decimal places, and \SI{88.9}{\percent} of NLOS cells have a
slope that does not move at all. Beyond \SI{80}{\metre}, where the diffraction
argument applies, the median slope goes from \num{2.994} to \num{2.992}. The
\SI{11.1}{\percent} of cells that do move all move downward, by up to
\num{1.0}, and they sit in the middle zone rather than the far one; they are
enough to pull the all-NLOS median from \num{2.99} to \num{2.76} while
leaving the far-zone median untouched.

With specular reflection and a single diffraction per path, nothing else in
the model can carry energy into the deep shadow, and diffraction weakens as
frequency rises. The far-zone discrepancy is therefore a property of the
propagation mechanisms included, not of how far they are followed.

This points at what is missing. A measured channel has a mechanism filling the
deep shadow that this simulation does not, and diffuse scattering from building
surfaces is the obvious candidate.

% =====================================================================
\section{Effect of Diffuse Scattering on NLOS Frequency Scaling}
\label{sec:scattering}

\subsection{Setting the scattering coefficient from surface roughness}

Diffuse scattering has a standard treatment in urban propagation prediction,
in which each surface is given a scattering coefficient $S$ and that
coefficient is fitted to measurement
\cite{degliesposti2001diffuse, degliesposti2004advanced}; measured values for
real building surfaces are reported in \cite{degliesposti2007measurement}.
Sionna RT 2.0.1 takes its material permittivity and conductivity from
ITU-R P.2040-3 \cite{iturp2040}, but that recommendation says nothing about
scattering coefficients; the value of zero that every material carries is a
choice in the implementation, set in \texttt{radio\_\hspace{0pt}material.py}.
Enabling diffuse reflection alone therefore changes nothing: at a fixed seed the map
moves by \SI{7e-6}{\decibel}, which is floating-point noise. The scattering
coefficient $S$ must be set explicitly, and it is therefore a modelling choice
rather than a switch. $S$ is the fraction of the reflected \emph{field} that
becomes diffuse, with the specular component scaled by $\sqrt{1-S^2}$, so
raising $S$ redistributes energy rather than adding it.

Fitting $S$ to measurement is not available here, and fitting it to 38.901
would make the comparison circular, so we derive it from surface roughness
instead. Under the Rayleigh criterion with the Ament attenuation factor
\cite{ament1953rough}, a surface of RMS height deviation
$\sigma_h$ attenuates the specular lobe by
\begin{equation}
\rho_s = \exp\!\left(-\frac{8\pi^2 \sigma_h^2 \cos^2\theta_i}{\lambda^2}\right),
\label{eq:ament}
\end{equation}
where $\theta_i$ is the angle of incidence measured from the surface normal,
$\lambda$ is the free-space wavelength, and $\sigma_h$ and $\lambda$ are in
the same length units. The energy leaving the specular lobe becomes diffuse,
giving $S = \sqrt{1-\rho_s^2}$. Because roughness is measured against
wavelength, a fixed surface implies a larger $S$ at higher carrier frequency.
Throughout the sweep we evaluate \eqref{eq:ament} at normal incidence,
$\theta_i = 0$, so $\cos^2\theta_i = 1$; at \SI{0.5}{\milli\metre} of
roughness this gives $S = 0.07$, \num{0.54} and \num{0.70} at \SI{3.5}{},
\SI{28}{} and \SI{39}{\giga\hertz}, where $\lambda$ is \SI{85.7}{},
\SI{10.7}{} and \SI{7.7}{\milli\metre}.

No single roughness value is defensible for a mixed urban scene, since the
roughness of real facade materials varies so widely. The comparison is
therefore swept across \SI{0.2}{} to \SI{2.0}{\milli\metre}. What that range
represents is not an arbitrary span but the transition itself: at the
wavelengths of interest here, it is the difference between a surface that
reflects specularly and one that scatters. The sweep is read in those terms
below rather than as an attempt to identify a correct value.
% TODO (Ali): read degliesposti2007measurement, then add one sentence here
% comparing the S values in Table IV against the scattering coefficients it
% measures on real building surfaces. A reviewer will ask whether the swept
% range is physically plausible, and this is the paper that answers it.
% Do not write that sentence until you have read the paper.

\subsection{Effect on the comparison with UMi}

Enabling diffuse scattering makes individual cells disagree far more between
runs than the bulk of the map does: at $S = 0.6$ the interquartile range of
the run-to-run difference is about \SI{1}{\decibel} while a few percent of
deep-shadow cells differ by up to \SI{44}{\decibel}. Three runs are too few
to characterise the tail of that distribution, so we describe the observed
spread rather than infer a shape from it, and each map below
is the per-cell median of three independent runs at \num{1e8} rays and every
reported quantity is a median or a share of cells rather than a mean or an
RMSE.

\begin{table*}[t]
\caption{Effect of roughness-derived diffuse scattering on agreement with
38.901 UMa. $S$ is implied by \eqref{eq:ament} at normal incidence. Interval
coverage uses UMa's NLOS shadow fading of \SI{6.0}{\decibel}. Every row,
including the scattering-off row, is the per-cell median of three independent
runs at \num{1e8} rays on one common mask; the off row agrees with the
independently traced \num{2e8}-ray map of Table~\ref{tab:umi} to
\SI{0.01}{\decibel}. The frequency slope involves no reference model.}
\label{tab:roughness}
\centering
\small
\begin{tabular}{lccccc}
\toprule
Roughness & $S$ at 3.5 / 28 / \SI{39}{\giga\hertz} & NLOS mean bias vs UMa (dB) & NLOS within UMa $\pm 1\sigma$ & \multicolumn{2}{c}{NLOS frequency slope} \\
\cmidrule(lr){5-6}
$\sigma_h$ &                                        & 3.5 / 28 / \SI{39}{\giga\hertz} & 3.5 / 28 / \SI{39}{\giga\hertz} & all & far ($>\SI{80}{\metre}$) \\
\midrule
off                    & 0 / 0 / 0            & $+9.4$ / $+14.7$ / $+15.6$ & \SI{15}{\percent} / \SI{15}{\percent} / \SI{15}{\percent} & 2.99 & 2.99 \\
\SI{0.2}{\milli\metre} & 0.03 / 0.23 / 0.32   & $+9.4$ / $+12.0$ / $+10.9$ & \SI{15}{\percent} / \SI{17}{\percent} / \SI{19}{\percent} & 2.27 & 2.70 \\
\SI{0.5}{\milli\metre} & 0.07 / 0.54 / 0.70   & $+9.4$ / $+7.6$ / $+6.5$   & \SI{15}{\percent} / \SI{31}{\percent} / \SI{41}{\percent} & 1.86 & 2.01 \\
\SI{1.0}{\milli\metre} & 0.15 / 0.86 / 0.96   & $+9.0$ / $+5.7$ / $+5.8$   & \SI{15}{\percent} / \SI{51}{\percent} / \SI{52}{\percent} & 1.76 & 1.88 \\
\SI{2.0}{\milli\metre} & 0.29 / 1.00 / 1.00   & $+7.7$ / $+5.7$ / $+5.9$   & \SI{19}{\percent} / \SI{52}{\percent} / \SI{52}{\percent} & 1.85 & 1.96 \\
\bottomrule
\end{tabular}
\end{table*}

Table~\ref{tab:roughness} shows the discrepancy of Sections~\ref{sec:umi}
and~\ref{sec:freq} closing as roughness rises, and closing in a frequency
selective way that follows directly from \eqref{eq:ament}.

Fig.~\ref{fig:roughness}(a) and~(b) plot the bias and the interval coverage
across the sweep. At \SI{3.5}{\giga\hertz} the implied $S$ never exceeds
\num{0.3} across the whole sweep, and the median far NLOS cell moves by less than
\SI{0.1}{\decibel}. The \SI{3.5}{\giga\hertz} bias is therefore essentially
unchanged, as it should be: at a wavelength of \SI{86}{\milli\metre}, surfaces
of this roughness are smooth. At \SI{28}{} and \SI{39}{\giga\hertz}, where the
wavelength falls to \SI{10.7}{} and \SI{7.7}{\milli\metre}, the same surfaces
are rough, and from \SI{0.5}{\milli\metre} upward diffuse paths add \SI{6}{} to
\SI{11}{\decibel} to the median far NLOS cell.

The consequence for the frequency slope is the result of this section, plotted
in Fig.~\ref{fig:roughness}(c), and it has to be read on the far-zone cells rather than on all NLOS cells, because
those are the ones the diffraction argument of Section~\ref{sec:freq} is
about. With scattering off, the far-zone median sits at \num{2.99}, the
single-edge value. At \SI{0.2}{\milli\metre} it has barely moved, to
\num{2.70}. It is only at \SI{0.5}{\milli\metre} that it drops to
\num{2.01}, and it then overshoots slightly to \num{1.88} and \num{1.96}.
There is no roughness in the swept range that places the far zone on the
UMi value of \num{2.13}: between \SI{0.2}{} and \SI{0.5}{\milli\metre} it
steps from above it to below it. It crosses the macrocell reference of
\num{2.00} in the same interval, reaching \num{2.01} at
\SI{0.5}{\milli\metre}, but Section~\ref{sec:scatlimits} shows that this
figure is not converged, so the crossing point should be read as bracketed
rather than located. Taken over all NLOS cells the same curve reads
differently, \num{2.27} at \SI{0.2}{\milli\metre} and \num{1.86} at
\SI{0.5}{\milli\metre}, because the near-zone reflection-lit cells sit near
\num{2.0} throughout and pull the median down.

Interval coverage follows a different trend from the slope, and the two should
not be read as one result. The share of NLOS cells inside UMa's one-sigma band
rises monotonically with roughness, from \SI{15}{\percent} with scattering
off to \SI{17}{} and \SI{19}{\percent} at \SI{0.2}{\milli\metre}, and to
\SI{51}{} and \SI{52}{\percent} by \SI{1.0}{\milli\metre}, where it
saturates. It does not reach the \SI{68}{\percent} the model implies, and
UMa's interval is the harder test: its NLOS shadow fading is
\SI{6.0}{\decibel} against UMi's \SI{7.82}{\decibel}, so the same maps score
\SI{61}{} and \SI{65}{\percent} against the microcell model. Either way the
roughness that brings the far-zone slope closest to the reference,
\SI{0.5}{\milli\metre}, is not the roughness that puts the most cells inside
the reference's own interval. A single roughness does not reproduce the model
on both measures at once.

The mechanism is therefore sufficient to explain the gap. A ray tracer without
diffuse scattering lights the deep shadow with a single diffracting edge and
inherits that edge's $30\log_{10}(f)$ scaling; a measured channel does not,
because real surfaces scatter, and scattering grows with frequency exactly
where diffraction weakens.

\begin{figure*}[!t]
\centering
\includegraphics[width=\textwidth]{fig_roughness_uma_ieee.png}
\caption{Roughness-derived diffuse scattering against the 38.901 UMa
comparison, as RMS surface roughness rises from off to
\SI{2.0}{\milli\metre}. (a) NLOS bias against UMa, per band. (b) Share of NLOS
cells inside UMa's own $\pm 1\sigma$ band, which is \SI{6.0}{\decibel} wide;
the model itself implies \SI{68}{\percent}. (c) Per-cell NLOS frequency
coefficient, median and 10th to 90th percentile, against the single-edge value
of \num{3.0} and UMa's \num{2.00}, which coincides with free space. (d) How
far (a) to (c) can be trusted: run-to-run disagreement in the far zone, and
how far the far-zone median still moves between \num{2e7} and \num{1e8}
rays.}
\label{fig:roughness}
\end{figure*}

\subsection{Ray budget and surface assumptions}
\label{sec:scatlimits}

Three limitations apply specifically to the scattering results.

The first is numerical, and it bounds how precisely the slope can be quoted;
Fig.~\ref{fig:roughness}(d) carries the two quantities it rests on, the
run-to-run disagreement in the far zone and the change in far-zone median
power between budgets.
With scattering off, the per-cell frequency slope is already converged: the
far-zone median is \num{2.99} both at \num{2e7} rays and at \num{1e8}, a
change of under \num{0.01}. With scattering on it is not. Measured on one
common mask, the far-zone median slope falls by \num{0.29}, \num{0.62},
\num{0.62} and \num{0.63} between those two budgets at \SI{0.2}{},
\SI{0.5}{}, \SI{1.0}{} and \SI{2.0}{\milli\metre} respectively. The sign is
the same at every roughness, so the scattering-on slopes reported above are
upper bounds and the converged values lie below them. The magnitudes behave
the same way: two runs differ in the far zone by a median of \SI{3.4}{} to
\SI{3.9}{\decibel}, and the median far cell rises by \SI{1.9}{} to
\SI{3.4}{\decibel} over the same budget range. What the sweep establishes is
therefore the direction and the approximate scale of the effect, not the
roughness at which the slope crosses the reference.

The second is the incidence assumption. Normal incidence in \eqref{eq:ament}
gives the largest $S$ a given roughness can imply, while canyon walls are
mostly struck obliquely, which would lower it, so the roughness needed for a
given effect is if anything higher than quoted.

The third is that a single $S$ is applied to every surface including the
ground, with Sionna's default Lambertian scattering pattern.

\section{Conclusion}
\label{sec:conclusion}

We have compared a ray-traced urban street canyon against the 38.901 urban
macrocell and microcell models at three carriers, cell by cell on one fixed
mask. Line of sight reproduces the macrocell model to within
\SI{0.3}{\decibel}; non line of sight does not, and the disagreement is
principally about frequency. Beyond \SI{80}{\metre}, \SI{81}{\percent} of
traced cells scale within \num{0.1} of $30\log_{10}(f)$, the signature of
single-edge diffraction, where the macrocell reference assumes
$20\log_{10}(f)$ and the microcell reference $21.3\log_{10}(f)$.
Raising the interaction depth does not change this, so it follows from which
propagation mechanisms are enabled rather than how far they are traced. Giving
every surface a scattering coefficient derived from its roughness has moved
the far-zone slope from \num{2.99} to \num{2.01} between \SI{0.2}{} and
\SI{0.5}{\milli\metre}, stepping past the reference value rather than
settling on it, while interval coverage has improved separately and
monotonically, saturating near \SI{52}{\percent} by \SI{1.0}{\milli\metre}. The
far-zone magnitudes are not fully converged, so these figures are indicative
rather than exact. Whether the $30\log_{10}(f)$ signature survives across
scenes, transmitter heights and street geometries is the obvious next
question.

% FALLBACK, for the single-author version if Thompson and Qi decline
% co-authorship. Uncomment this and cut their blocks from \author above, and
% change "authors" to "author" in the two statements below.
%
% \section*{Acknowledgements}
% The author thanks Prof.\ John S.\ Thompson and Tianrun Qi of the Institute for
% Imaging, Data and Communications, University of Edinburgh, for proposing the
% cross-frequency comparison and the diffuse scattering study, and for their
% guidance throughout.

\section*{Funding}
This work received no specific grant from any funding agency, and was carried
out on freely available software and ordinary desktop hardware.

\section*{Conflict of Interest}
The authors declare that they have no conflicts of interest.

\section*{Data and Code Availability}
All code, scene configurations, raw sweep outputs and known-answer checks
needed to reproduce every figure and table in this paper are publicly
available at
\url{https://github.com/alioramuss/radiomap-sparse-reconstruction}
% TODO (Ali): tag a release before submission and put its tag and commit
% hash here. Qi asked for a citable version rather than a bare repository
% link, and he is right that a moving main branch is not a citation.
(release and commit to be tagged at submission). Table~\ref{tab:convergence}
is produced by \texttt{scripts/06\_\hspace{0pt}multiband\_\hspace{0pt}convergence.py},
Tables~\ref{tab:umi} and~\ref{tab:freqcoef} by
\texttt{scripts/14\_\hspace{0pt}umi\_\hspace{0pt}comparison.py},
Table~\ref{tab:depth} by \texttt{scripts/17\_\hspace{0pt}depth\_\hspace{0pt}check.py},
and Table~\ref{tab:roughness} by
\texttt{scripts/18\_\hspace{0pt}roughness\_\hspace{0pt}sweep.py} and
\texttt{scripts/19\_\hspace{0pt}roughness\_\hspace{0pt}analysis.py} with
\texttt{-{}-model uma}.

\bibliographystyle{IEEEtran}
\bibliography{refs}

\end{document}
